\chapter{Validation} % Chapter 5 (renumbered)
\label{Validation}

% =============================================================================
% VALIDATION (5-8 pages)
% Validates reasoning trace integrity, K-level reasoning manipulation,
% and addresses faithfulness concerns. This chapter ensures the thesis
% can defend its use of reasoning traces as behavioral data.
% =============================================================================

% -----------------------------------------------------------------------------
% SECTION 6.1: Reasoning Trace Analysis
% -----------------------------------------------------------------------------
\section{Reasoning Trace Analysis}

% TODO: 6.1.1 — Trace content overview
\subsection{Trace Content by Reasoning Level}
%   - Token count distributions: L0 vs L1 vs L2 vs L3
%   - Qualitative examples of trace content per level
%   - Does L2 actually mention opponent modeling? Does L3 show recursion?
%   - Manipulation check: do higher levels produce qualitatively different traces?

% TODO: 6.1.2 — Trace coding methodology
\subsection{Systematic Trace Coding}
%   - LACA framework (Chew et al., 2023): codebook → LLM coding → human verification
%   - Codebook categories: payoff calculation, opponent prediction, recursive reasoning,
%     threat assessment, cooperation signaling, resource tracking
%   - Inter-rater reliability with Gwet's AC1
%   - Frequency of each category per reasoning level
%   - Key question: does L2 contain MORE opponent modeling than L1?
%     Does L3 contain MORE recursive reasoning than L2?

% -----------------------------------------------------------------------------
% SECTION 6.2: Faithfulness Validation
% -----------------------------------------------------------------------------
\section{Faithfulness of Reasoning Traces}

% TODO: 6.2.1 — The faithfulness problem
\subsection{The Faithfulness Problem}
%   - CoT traces may not reflect internal computation (Turpin et al., 2023)
%   - Reasoning models construct post-hoc rationales (Chen et al., 2025)
%   - CoT ≠ mechanistic explanation (Barez et al., 2025)
%   - OUR FRAMING: traces are behavioral data, not explanations
%     We measure what agents DO when given different computational depth,
%     not what they "really think"

% TODO: 6.2.2 — Truncation test (Lanham et al., 2023)
\subsection{Truncation Test}
%   - Method: for a subset of runs, truncate the thinking trace at various points
%     and check if the final action changes
%   - If truncating the opponent-modeling section of an L2 trace changes the action,
%     that section has causal influence on the output
%   - Expected: L0 traces trivially faithful (short, heuristic)
%     L2/L3 traces: some sections causally relevant, others rationalizing
%   - Report: fraction of traces where truncation changes action, per level

% TODO: 6.2.3 — Thought Anchors (Bogdan et al., 2025)
\subsection{Thought Anchors Analysis}
%   - Method: counterfactual resampling at sentence level
%   - Identify "thought anchors": sentences that, when resampled, change the outcome
%   - Apply to subset of critical rounds (e.g., first attack, cooperation breakdown)
%   - Expected: planning/uncertainty sentences are anchors, not filler
%   - This is a stronger test than truncation (Bogdan et al., 2025)

% TODO: 6.2.4 — TextGrad-inspired analysis
\subsection{Gradient-Based Prompt Sensitivity}
%   - Method: find minimal prompt perturbations that flip the action
%   - If small changes to the reasoning instruction change behavior,
%     the instruction has genuine causal influence
%   - Report: sensitivity per reasoning level
%   - Limitation: computational cost, may only be feasible on subset

% -----------------------------------------------------------------------------
% SECTION 6.3: K-Level Reasoning Verification
% -----------------------------------------------------------------------------
\section{K-Level Reasoning Verification}

% TODO: 6.3.1 — Does the manipulation work?
\subsection{Manipulation Check}
%   - L0 should NOT predict opponents (instruction says "do not deliberate")
%   - L1 should compute payoffs but NOT model opponents
%   - L2 should predict opponent actions
%   - L3 should reason about what opponents think about YOU
%   - Check via trace coding: are these instructions followed?
%   - Report: compliance rate per level

% TODO: 6.3.2 — K-Level depth measurement
\subsection{Depth Measurement}
%   - Following Zhang et al. (NAACL 2025):
%     Count levels of recursive reasoning in each trace
%   - L0 should show depth 0, L2 should show depth 1, L3 should show depth 2
%   - Do agents actually recurse to the instructed depth?
%   - Or do they default to a lower level despite the instruction?

% TODO: 6.3.3 — Compliance vs behavior
\subsection{Compliance vs Behavioral Outcomes}
%   - Even if L3 agents don't fully follow recursive instructions,
%     they still BEHAVE differently from L0/L1
%   - The behavioral effect is what matters for the thesis
%   - The trace compliance is a secondary validation
%   - Frame: "the prompt manipulation produces reliably different behavior
%     and reasoning patterns, even if full recursive compliance is partial"

% -----------------------------------------------------------------------------
% SECTION 6.4: Robustness and Sensitivity
% -----------------------------------------------------------------------------
\section{Robustness Checks}

% TODO: 6.4.1 — Prompt sensitivity
\subsection{Prompt Sensitivity}
%   - Test semantically equivalent reformulations of L0-L3 instructions
%   - Report FormatSpread or PromptSensiScore
%   - Key question: are results robust to wording, or fragile?

% TODO: 6.4.2 — Seed stability
\subsection{Seed Stability}
%   - Across-run variance at same condition: ICC analysis
%   - How many runs needed for stable effect estimates?
%   - Report: variance decomposition (within-run vs between-run)

% TODO: 6.4.3 — Generative ABM validation
\subsection{Generative ABM Validity}
%   - Following Larooij et al. (2025): focus on macro-level patterns
%   - Do aggregate metrics (Gini, cooperation rate) show consistent trends
%     across runs, even if individual agent behavior varies?
%   - Are phase transitions reproducible across seeds?
